\section{Introduction}
The space weather environment is home to many hazards that put humans and their artificial satellites at risk \cite{eastwood2017economic}. 
Solar Energetic Particles (SEPs) are bursts of high energy protons $>10$ MeV, electrons and $\alpha$ particles following the magnetic field lines from the Sun reaching relativistic speeds $(\approx 0.9c)$. Coronal Mass Ejections (CMEs) drive shock waves that accelerate SEPs to tens of GeV and flux $\geq 10$ MeV making them difficult to shield \cite{reames2021solar}.
SEPs are one of many species in the extra-terrestrial ecosystem; however, protons $>10$ MeV are one of the greatest offenders damaging electronic components on satellites, aircraft, high frequency communications, and human health, making reliable forecasting essential \cite{eastwood2017economic}.
SEP relativistic  electrons ($0.95c$, $1$ MeV) are always present during storm events and precede slower incoming SEP ($0.25c$, $30$ MeV) protons due to less scattering  \cite{nunez2018hesperia}.
This simple relationship has been the basis of a portion of forecast models predicting high energy protons at $10$ pfu.

Physics models that simulate heliospheric transport and particle acceleration processes are computationally expensive and are usually studied to understand underlying processes rather than generating predictions \cite{Camporeale2019,kasapis2026review}.
The issue with all physics-based SEP models is that they may have poorly constrained free parameters such as diffusion coefficients and which specific magnetic field line to select to produce a prediction for the particle spectrum at the observer's location \cite{whitman2023}.
Empirical and statistical models provide predictions but rely heavily on historical correlations leaving probabilities of an event to occur \cite{whitman2023}. Together, physics, statistical, and empirical models are limited by uncertain inputs and must trade off realistic behavior with forecast timeliness \cite{yu2026}.

Machine Learning (ML) models have had major breakthroughs in the application of space weather due to the large availability of in situ observations and recent advancements in Graphics Processing Unit (GPU) computing \cite{Camporeale2019}. ML models predicting SEPs with space weather properties demonstrated improved results establishing reliable forecast methods \cite{Posner2024,yu2026,nunez2022umasep,Dayeh2024}.
With the recent development of self-attention architectures, transformers, \cite{vaswani2017attention}, there is still room for improvements. Recurrent models struggle with long-range dependencies, convolution models require deep stacking to enlarge receptive fields, and both scale suboptimally to high-dimensional and high-volume data \cite{zhao2026transformersurvey}. Transformers evaluate the totality of input data in parallel to generate reliable, accurate outputs. Their architecture allows for optimal scaling \cite{vaswani2017attention}. Because the introduction of transformers is recent, there are limited applications in predicting space weather, therefore a fully realized comparative analysis of forecasting performance is necessary and would provide much insight for other modelers.

We analyze   time series  SEP data (5-minute cadence) recorded by the Solar and Heliospheric Observatory Comprehensive Suprathermal and Energetic Particle Analyzer--Electron Proton Helium Instrument
(SOHO-COSTEP-EPHIN) located at Lagrange Point 1 (L1). We compare transformer models to a collection of SEP forecasting models of all types under a battery of assessments. These include multilayer perceptron neural network (NN) and recurrent neural network (RNN) models \cite{torres2025machine}, a physics model using an electron increase parameter matrix \cite{posner2007forecasting}, and simple statistical and empirical models, persistence and linear regression.  
Our tests include a fixed 2-hour time window of particle fluxes as an input to predict a single proton flux value at either $30$ minutes or $60$ minutes ahead. Additionally, 
expanding the input time window, increasing forecasting time, and parameter sizes of our models will give us a holistic diagnosis of the performances
between the models.

Across  all forecasting models types, each model reads in solar storm parameters, measures resulting metrics, and performs analysis that is different than the next. No two model studies
analysis is the same. Some studies do not include a naive persistence model to compare against. This study also hopes to establish standards of metric acquisition and scoring
to properly compare forecasting models to each other similar to the time series transformer field. 

Section \ref{sec:relatedwork} will discuss related work on transformer applications.
Section \ref{sec:approach}  will report which transformer models we developed and why,  and reintroduce previously mentioned models more in depth. A summary of selected metrics will be provided to compare forecasting accuracy. Finally, we provide an outline of our battery assessment in detail.
Section \ref{sec:experimentalevaluation} will show our evaluation of the model performances for each test with analysis.
Section \ref{sec:conclusion} will summarize and discuss future work, potential improvements and limitations.
